6. (1 pt) Determine las condiciones de \( a, b, c \in \mathbb{Z} \) para que
\[
\operatorname{det}\left(\begin{array}{ccc}
a+b & c & c \\
a & b+c & a \\
b & b & a+c
\end{array}\right)
\]
sea múltiplo de 8 . \\

\begin{solution}
Para determinar las condiciones sobre \( a, b, c \in \mathbb{Z} \) para que el determinante de la matriz dada sea múltiplo de 8, seguimos los siguientes pasos:
Consideremos la matriz:
\[
A = \begin{pmatrix}
a+b & c & c \\
a & b+c & a \\
b & b & a+c
\end{pmatrix}
\]
Calculamos el determinante de A:
\[
\det(A) = \begin{vmatrix}
a+b & c & c \\
a & b+c & a \\
b & b & a+c
\end{vmatrix}
\]
Expandimos el determinante utilizando la regla de Sarrus:
\[
\det(A) = (a+b)\left( (b+c)(a+c) - a \cdot b \right) - c\left( a(a+c) - a \cdot b \right) + c\left( a \cdot b - b(b+c) \right)
\]
Simplificamos cada término:

Calculamos (b+c)(a+c)−a⋅b:

\[
(b+c)(a+c) = ab + ac + bc + c^2
\]
\[
(ab + ac + bc + c^2) - ab = ac + bc + c^2
\]

Multiplicamos por (a+b):

\[
(a+b)(ac + bc + c^2) = a^2c + abc + ac^2 + abc + b^2c + bc^2
\]

Calculamos −c(a(a+c)−a⋅b):

\[
a(a+c) = a^2 + ac
\]
\[
a^2 + ac - ab = a^2 + ac - ab
\]
\[
-c(a^2 + ac - ab) = -a^2c - a c^2 + a b c
\]

Calculamos c(a⋅b−b(b+c)):

\[
b(b+c) = b^2 + bc
\]
\[
a b - (b^2 + bc) = a b - b^2 - bc
\]
\[
c(a b - b^2 - bc) = a b c - b^2 c - b c^2
\]
Sumamos todos los términos:
\[
\det(A) = (a^2c + 2abc + a c^2 + b^2c + b c^2) + (-a^2c - a c^2 + a b c) + (a b c - b^2 c - b c^2)
\]
Simplificando:
\[
\det(A) = 4 a b c
\]
Queremos que det(A) sea múltiplo de 8:
\[
4 a b c \equiv 0 \pmod{8}
\]
Dividimos ambos lados por 4:
\[
a b c \equiv 0 \pmod{2}
\]
Esto implica que el producto abc debe ser par. Es decir, al menos uno de los números a,b o c debe ser par.
\end{solution}